\subsection{The boundary fibre in the asserted Bers inverse limit}

\begin{result}{E.48}{Refuted; exact boundary-fibre replacement proved}{The discrete diagonal model loses a completed stabilizer}
Let $\Gamma$ be a stable mapping class group with a nonempty
boundary stratum, let $X$ be its Bers bordification, and put
$G=\check\Gamma$.  Let $N\triangleleft_oG$ range over open
normal subgroups and set $\Gamma_N=\Gamma\cap N$.
There is a natural map of sets
\[
 \Phi:(X\times G)/\Gamma\longrightarrow
                 \varprojlim_N X/\Gamma_N,
 \qquad
 \Phi_N([x,g])=[\gamma_N^{-1}x],
 \quad \gamma_N\equiv g\pmod N.                 \tag{E.248}
\]
The action in the source is the diagonal left action.
This map is not injective.  In particular, the natural
isomorphism with this diagonal quotient asserted in the
proof of Boggi's Theorem~6.3, PDF p.~29,
\cite{ET_Boggi2014} is false as a statement about these
underlying sets.

For a fixed boundary point $p\in X$, write
$S=\operatorname{Stab}_\Gamma(p)$.  The exact replacement
on its orbit is
\[
 \Gamma\backslash((\Gamma p)\times G)
       \cong G/S
       \longrightarrow G/\overline S
       \xrightarrow{\sim}
           \varprojlim_N\Gamma_N\backslash\Gamma/S.
                                                               \tag{E.249}
\]
Here the middle map is the actual quotient map and
$\overline S$ is closure in $G$.  Formula (E.249) is a
statement about the fixed orbit and its finite-level
completion.  No global corrected Bers model is asserted.
\end{result}
\begin{proof}
Density of $\Gamma$ in $G$ supplies $\gamma_N$ in
(E.248).  Replacing it by another lift changes
$\gamma_N^{-1}x$ by an element of the normal subgroup
$\Gamma_N$.  Replacing $(x,g)$ by $(\eta x,\eta g)$,
with $\eta\in\Gamma$, permits the lift $\eta\gamma_N$
and leaves $\gamma_N^{-1}x$ unchanged.  Thus the map is
well defined; using the same lifts modulo smaller
quotients proves compatibility in the inverse system.

Choose a node $\gamma$ of $p$.  Its integral Dehn twist
$t_\gamma$ fixes $p$, since twisting the collar which
has been pinched changes no marked nodal surface.
The valid twist-lattice part of the cutting theorem,
Boggi Theorem~4.5, gives an embedding
\[
 \Delta_\gamma=\widehat{\mathbb Z}\,t_\gamma
                         \hookrightarrow G.               \tag{E.250}
\]
This group is uncountable, whereas $\Gamma$ is countable.
Choose $h\in\Delta_\gamma\setminus\Gamma$.
The two classes $[p,1]$ and $[p,h]$ in the source of
(E.248) are distinct: equality would require some
$\eta\in\Gamma$ with $\eta p=p$ and $\eta=h$.
For each $N$, density of the integral twists in
$\Delta_\gamma$ supplies an integer $a_N$ such that
$h\equiv t_\gamma^{a_N}\pmod N$.  Choose this twist
as $\gamma_N$ in (E.248).  It fixes $p$, so
\[
       \Phi_N([p,h])=[p]=\Phi_N([p,1])
                       \qquad\text{for every }N.          \tag{E.251}
\]
This proves noninjectivity without any assertion about
$\Gamma\cap\Delta_\gamma$, and without using the later
twist-power results whose proof is under examination.

For completeness, identify the orbit comparison exactly.
The map $gS\mapsto[p,g^{-1}]$ identifies $G/S$ with
the first term of (E.249).  It is onto because every
first coordinate lies in $\Gamma p$, and its equality
relation is exactly right multiplication by $S$.
At level $N$, the orbit of $p$ is
$\Gamma_N\backslash\Gamma/S$, identified with
$(G/N)/q_N(S)$.  Since this target is finite,
$q_N(S)=q_N(\overline S)$.  The map
\[
 G/\overline S\longrightarrow
          \varprojlim_N (G/N)/q_N(S)                       \tag{E.252}
\]
is injective: equal images mean that the difference
of two representatives lies in every $N\overline S$,
whose intersection is $\overline S$.  It is onto by
compactness: a compatible family of cosets specifies
a family of nonempty closed subsets of $G$ with the
finite-intersection property, because the indexing
system admits common refinements.  Their intersection
gives a representative of the family.  Compactness
and the Hausdorff property make (E.252) a homeomorphism
for the profinite quotient topology.  With the source
identification just given, (E.248) restricts to the
quotient map $G/S\to G/\overline S$.

The precise remaining source obligation in E.45 is
therefore not settled by that printed proof.  For
primitive twists, E.47 proves
$C_G(t_\alpha)=\operatorname{Stab}_G(\alpha)$.
To deduce the commuting criterion from this equality
one still needs to prove
\[
 t_\beta\in\operatorname{Stab}_G(\alpha)
   \quad\Longrightarrow\quad
 \beta=\alpha\ \text{or}\ \{\alpha,\beta\}
                   \text{ is a profinite simplex}.         \tag{E.253}
\]
The counterexample above is to the asserted inverse-limit
identification used in that derivation; it is not a
counterexample to (E.253).  The positive multitwist
uniqueness theorem in Boggi--Zalesskii
\cite[Theorem~6.1]{ET_BZ2019Magnus} assumes that both
supports already are multicurves.  Its Corollary~6.2
then cites the same Boggi Corollary~6.4 and Theorem~6.6.
Consequently those displayed comparisons do not by
themselves supply the missing implication (E.253).
\end{proof}
